% !TeX program = pdflatex
%%% doc/bangor.tex — documentation for the bangor class, typeset with the class
\documentclass[dissertation,tables,glossary,12pt,a4paper]{bangor}

\title{The bangor class}
\author{Bangor University LaTeX templates project}
\date{2026/09/24}
\shorttitle{The bangor class}

\addabbrev{PDF}{portable document format}
\addabbrev{TDS}{TeX directory structure}

\begin{document}

\begin{titlepage}
  \centering
  \bangoridentityblock[3.5cm]
  \vfill
  {\LARGE\bfseries The bangor class\par}
  \vfill
  {\large Bangor University LaTeX templates project\par}
  \vfill
\end{titlepage}
\pagenumbering{roman}

\tableofcontents
\printabbrevs

\chapter{What this class does}
The \textsf{bangor} class typesets documents for Bangor University. It is one file. It gives you the brand, the document identity, the regulation layout, the statutory declarations, table support, an abbreviation list, and a headed letter style. This document is written with the class itself, so every example that it shows also runs here. It is written in controlled English, so that readers who do not use English as a first language can read it without difficulty.

The class has two modes. A degree option, for example \verb|phd|, selects the thesis mode. The option \verb|letter| selects the letter mode. Do not mix the two: a thesis option in a letter is an error that names the option.

\section{Install}
Put \verb|bangor.cls| and \verb|bangor-crest-colour.pdf| in the folder of your document. That is all. The class needs no other file from this project. If you install the class in a TeX system, keep the two files together, because the class looks for the crest beside itself.

If the crest file is missing, the class does not stop. It prints a frame that holds the wordmark where the crest belongs, and it gives one warning that names the download address.

\section{The two modes}
For a thesis or dissertation, write the class line with a degree option:

\begin{quote}
  \verb|\documentclass[phd,12pt,a4paper]{bangor}|
\end{quote}

For a headed letter, write the class line with the letter option:

\begin{quote}
  \verb|\documentclass[letter,a4paper,11pt]{bangor}|
\end{quote}

The two optional features are also class options. Write \verb|tables| for the \verb|longtabular| environment and \verb|glossary| for the abbreviation list. This document uses both.

\chapter{The brand layer}
The brand layer belongs to both modes. It gives you three things.

\begin{itemize}
  \item The colours \verb|bangorred| and \verb|bangorsun|, from the university brand hub.
  \item The crest: \verb|\bangorcrest[3cm]|. The optional argument sets the width. The default width is 3cm.
  \item The wordmark: \verb|\bangorwordmark|. This prints the university name in bold, in black by default. Pass a colour name for an accent.
\end{itemize}

The version stamp macro \verb|\bangorversionstamp| records the class version. The class places this stamp in the PDF metadata, so every generated \useabbrev{PDF} records what produced it.

\chapter{Document identity}
Set the document facts in the preamble. School and college are free text, because the university renames its schools and colleges.

\begin{itemize}
  \item \verb|\school{School of Computer Science and Engineering}|
  \item \verb|\college{College of Science and Engineering}|
  \item \verb|\degreescheme{Doctor of Philosophy}|
  \item \verb|\supervisor{Dr A. Supervisor}|. Call this command once for each supervisor.
  \item \verb|\submissiondate{September 2026}|
  \item \verb|\submissionwordcount{9930}|. Use this for taught work that must state the word count.
  \item \verb|\shorttitle{Short title}| sets the title that prints in the footer. If you do not set it, the footer shows the full title.
\end{itemize}

The command \verb|\bangoridentityblock| renders the crest, the wordmark, and the school and college lines. The thesis title page and the letter letterhead both use this block, so neither repeats the brand facts.

\chapter{Regulation layout}
The class applies the rules from Regulation 03, 2025 Version 01, section 6.3, and checks them.

\begin{itemize}
  \item Fonts: \verb|font=termes| (Times-Roman equivalent, the default), \verb|font=bonum| (Bookman equivalent), or \verb|font=heros| (Arial and Helvetica equivalent, sans-serif). Serif body text must be at least 12pt. Sans-serif body text must be at least 10pt.
  \item Margins: 2.5cm at the binding edge, 2cm on the other sides.
  \item Line spacing: 1.5 in the main text. Set the abstract, indented quotations, and footnotes single spaced if you want; the regulation permits this.
\end{itemize}

Strict mode is on by default. A rule that the document does not meet fails the build. Set \verb|strict=false| to turn each failure into a warning. The class checks the settings that it controls. It cannot check changes that you make around it.

\chapter{Statutory declarations}
The command \verb|\statementspage| prints the Statement of Originality and the Statement of Availability. The originality text reproduces the statutory declaration from the regulation exactly. Use the \verb|welsh| class option to print the Welsh text as well.

The generative AI statement is a fill-in template. Replace it with your own statement:

\begin{quote}
  \verb|\aistatement{Your statement about your use of AI tools.}|
\end{quote}

This is the Statement of Originality as this class prints it:

\originalitystatement

\chapter{Tables}
Use the class option \verb|tables|. It provides the \verb|longtabular| environment and loads the packages behind it: \textsf{xltabular} and \textsf{pdflscape}.

\verb|longtabular| is one environment for a table that may run past a page, need \verb|X| columns to fit the text width, or both. It paginates at row boundaries, repeats the header on each continuation page when you place the header before \verb|\endhead|, and sizes \verb|X| columns to the text width. The table below is a \verb|longtabular|.

\begin{longtabular}{l X r}
  \caption{A table that uses an X column.}\label{tab:xcolumn} \\
  \hline
  Item & Description & Value \\
  \hline
  \endfirsthead
  \hline
  Item & Description & Value \\
  \hline
  \endhead
  Margin & The binding edge is at least 2.5cm and the other sides are at least 1.5cm. & 2.5 \\
  Spacing & The main text uses 1.5 line spacing. & 1.5 \\
  Size & Serif body text is at least 12pt. & 12 \\
\end{longtabular}

The optional argument turns on rotation for wide tables: \verb|\begin{longtabular}[wide]{l X r}|. It also takes \verb|columns=N|, which wraps a narrow table's rows into N side-by-side panes on one page, so a long narrow table uses the width of the page instead of running onto the next page. Each pane repeats the head, the caption prints once above the panes, and the wrapped table must fit one page. In this form the head sits between \verb|\panehead| and \verb|\endpanehead| in the body, because the longtable block markers cannot be read again outside a table. A missing marker is an error. The options combine as \verb|[wide,columns=2]|. Table~\ref{tab:wrapped} uses two panes.

\begin{longtabular}[columns=2]{l l r}
  \caption[Wrapped inventory]{An inventory wrapped into two columns on one page.}\label{tab:wrapped}\\
  \panehead
  \hline
  Item & Kind & Value \\
  \hline
  \endpanehead
  a01 & first kind & 1 \\
  a02 & second kind & 2 \\
  a03 & third kind & 3 \\
  a04 & fourth kind & 4 \\
  a05 & fifth kind & 5 \\
  a06 & sixth kind & 6 \\
\end{longtabular}

The underlying \verb|longtable|, \verb|tabularx| and \verb|landscape| environments remain available for special cases. The only permitted optional items are \verb|wide| and \verb|columns=N|. Any other item is an error.

\chapter{Abbreviations}
Use the class option \verb|glossary|. It loads no package, so it adds almost no compile time.

\begin{itemize}
  \item \verb|\addabbrev{ML}{machine learning}| registers an abbreviation.
  \item \verb|\useabbrev{ML}| prints the full form the first time and the short form after that.
  \item \verb|\Useabbrev{ML}| does the same with a capital letter on the first word. Use it at the start of a sentence.
  \item \verb|\printabbrevs| prints an unnumbered heading, entered in the contents only with \verb|frontmattertoc|, above a list of the abbreviations in alphabetical order. Each entry ends with the pages that use it, and a run of three or more consecutive arabic pages is written as a range. The heading reads ``List of Abbreviations'' (the text of \verb|\abbrevname|) unless an optional argument gives another, as in \verb|\printabbrevs[Glossary]|. The page lists come from the previous run, so a first run shows none.
\end{itemize}

The list at the front of this document is the result. It has two entries: the \useabbrev{TDS} layout that TeX systems use to install a class, and the \useabbrev{PDF} that this document produces.

The command \verb|\captiondesc[short]{title}{description}| is the caption form with a description paragraph beneath the title. The optional argument is the list-of-floats entry and defaults to the title. A list entry over 80 characters is a regulation violation: a build error in strict mode, a warning otherwise.

\chapter{Class options}
The degree keywords \verb|phd|, \verb|mphil|, \verb|mres|, \verb|mscres|, \verb|engd|, and \verb|profdoc| set the wording ``thesis''. The taught keywords, for example \verb|msc|, \verb|bsc|, and \verb|beng|, set the wording ``dissertation''. The options \verb|thesis| and \verb|dissertation| override the wording directly.

\begin{itemize}
  \item \verb|citations=style| selects the biblatex style. The default is \verb|ieee|. Any biblatex style name is valid. Ask your supervisor which style your school requires.
  \item \verb|font=termes|, \verb|strict|, and \verb|strict=false| set the font and the validation mode.
  \item \verb|print| adds a blank page before the title page, for duplex printing. \verb|digital| is the default and prints no blank page.
  \item \verb|welsh| prints the statutory declaration in Welsh and loads babel with Welsh.
  \item \verb|headers=chapter-section| is the default. The header shows the chapter name on the left and the section name on the right, and the page number is in the footer. \verb|headers=chapter| shows the chapter name and the page number in the header.
  \item \verb|headrule| and \verb|footrule| set the rule below the header and the rule above the footer. Each takes \verb|hairline| (the default), \verb|visible|, or \verb|none|. The two options work separately.
  \item \verb|pageofm| is on by default and prints the page number as ``Page N of M''. \verb|pageofm=false| prints the number alone.
  \item \verb|appendixpages| is on by default and numbers each appendix separately. The title page carries no number, the front matter is numbered in roman numerals, and the numbering restarts in arabic at the first numbered chapter. The count restarts at each appendix and the footer reads ``Appendix A - Page n of m'', where m is the number of pages in that appendix. Then M in ``Page N of M'' is the last page of the main text, not of the document. \verb|appendixpages=false| continues the arabic sequence through the appendices.
  \item \verb|sharedfloats| is on by default and puts figures, tables, listings, and algorithms on one counter, so the sequence reads Figure 1, Table 2, Listing 3. \verb|sharedfloats=false| gives each type its own counter. Listings and algorithms take part when their packages are loaded.
  \item \verb|sectionfloats| is on by default and numbers floats within their section, so the first float in section 5.1 is 5.1.1. A float before the first section of a chapter is numbered by the chapter alone (5.1). With \verb|sectionfloats=zero| it is numbered by section zero (5.0.1), and with \verb|sectionfloats=none| floats are numbered within their chapter. The two options combine: together they give one sequence that restarts at each section.
  \item \verb|frontmattertoc| enters the unnumbered front matter in the contents: the abstract, the acknowledgements, the list of abbreviations, and the lists of figures and tables. Without it the contents lists the numbered chapters, the bibliography, and the appendices, and each front matter heading still gets a PDF bookmark.
  \item All other options pass to the \textsf{report} class, or to the \textsf{letter} class in the letter mode.
\end{itemize}

\chapter{Thesis commands}
Set the metadata in the preamble, then call \verb|\bibliographysetup| once, add your resources with \verb|\addbibresource|, and build the front matter:

\begin{quote}
  \verb|\maketitle| \quad \verb|\statementspage| \quad \verb|\begin{acknowledgements}| \verb|...| \verb|\end{acknowledgements}| \quad \verb|\begin{abstract}| \verb|...| \verb|\end{abstract}| \quad \verb|\tables|
\end{quote}

The abstract is single spaced. The regulation caps it at 600 words.

\section{Headers and footers}
The footer shows the title. Set \verb|\shorttitle{Short title}| in the preamble to show a shorter title.

The header shows the chapter name and, in the default layout, the section name. A name that is too wide for the header stops the build. The error message names the chapter or the section. Give a shorter name with the optional argument: \verb|\chapter[Short title]{Long title}|. The short title prints in the header. It also prints in the table of contents, because the LaTeX kernel uses the optional argument for both.

A title that is too wide for the footer also stops the build. Set \verb|\shorttitle| to fix it. With \verb|strict=false| the build continues: the class cuts the text short and shows a warning.

To change the layout in other ways, use the \textsf{fancyhdr} commands in your preamble.

\chapter{Headed letters}
Use the class line with the \verb|letter| option, as in the chapter that starts this document. The class loads the standard \textsf{letter} class for you. Put \verb|\bangorletterhead| at the top of each letter, before \verb|\opening|. The command prints the crest, the wordmark, the school line if you set one, and a rule in the brand red. The Bangor letter template on GitHub is a complete example.

A letter takes the options of the \textsf{letter} class, for example \verb|a4paper| and \verb|11pt|. It takes no thesis option.

\chapter{Version stamping in generated PDFs}
The class places the version stamp in the PDF metadata, in the subject and the keywords fields. The stamp records the class version, and the template version when the release process provides one. Open the document information panel of the PDF to see which version produced it.

\end{document}
